\chapter{动能与势能}

\section{章节主线与学习目标}

本章研究力的空间累积效应与机械能转化规律。上一章的动量定理强调力对时间的累积：
\[
    \boldsymbol I=\int \boldsymbol F\,\mathrm dt,
\]
而本章的功与能量方法强调力对位移的累积：
\[
    W=\int \boldsymbol F\cdot\mathrm d\boldsymbol r .
\]
二者都来源于牛顿第二定律，但适用对象和物理意义不同。

\subsection{本章主线}

动能与势能一章可以按如下脉络学习：
\[
    \text{功与功率}
    \longrightarrow
    \text{动能定理}
    \longrightarrow
    \text{保守力与势能}
    \longrightarrow
    \text{功能原理}
    \longrightarrow
    \text{机械能守恒定律}.
\]
本节之后将先讨论功与功率，再由牛顿第二定律推出动能定理。保守力、势能、功能原理以及机械能守恒定律将在后续章节中继续展开。

\subsection{学习目标}

学习本章前半部分，应掌握以下内容：

\begin{enumerate}
    \item 理解功的定义
    \[
        W=\int_a^b\boldsymbol F\cdot\mathrm d\boldsymbol r,
    \]
    明确“只有力沿位移方向的分量做功”；
    \item 掌握瞬时功率
    \[
        P=\boldsymbol F\cdot\boldsymbol v,
    \]
    理解功率描述做功快慢，而不是力的大小；
    \item 掌握质点动能定理
    \[
        W_{\mathrm{net}}=\Delta E_k,
    \]
    并能用它处理已知位移、求速率的问题；
    \item 区分动量方法与能量方法：冲量改变动量，功改变动能；
    \item 能根据题目条件选择牛顿第二定律、动量定理或动能定理。
\end{enumerate}

\subsection{动量方法与能量方法的选择}

若题目强调力作用的时间、碰撞过程中的平均力或动量变化，通常优先考虑动量定理：
\[
    \boldsymbol I=\Delta\boldsymbol p.
\]
若题目强调运动路径、位移、速度大小变化或不需要求时间，通常优先考虑功与动能定理：
\[
    W_{\mathrm{net}}=\Delta E_k.
\]
简单地说，出现 $\Delta t$ 时先想到冲量，出现 $\mathrm d\boldsymbol r$ 或路径长度时先想到功。

\newpage

\section{功与功率}

功描述力在空间位移上的累积效应。不能把功简单理解成“力的大小乘以路程”，因为只有力在位移方向上的分量参与做功。功率则描述做功的快慢。

\subsection{恒力做功}

若物体在恒力 $\boldsymbol F$ 作用下发生位移 $\boldsymbol s$，力与位移夹角为 $\theta$，则该力做功为
\[
    W=\boldsymbol F\cdot\boldsymbol s=Fs\cos\theta .
\]
因此：

\begin{enumerate}
    \item 当 $0\leq \theta<90^\circ$ 时，$W>0$，力做正功；
    \item 当 $\theta=90^\circ$ 时，$W=0$，力不做功；
    \item 当 $90^\circ<\theta\leq 180^\circ$ 时，$W<0$，力做负功。
\end{enumerate}

功是标量。功的正负表示能量转移方向，不表示空间方向。

\subsection{变力做功}

若力随位置变化，或物体沿曲线运动，则应把路径分成许多微小位移元 $\mathrm d\boldsymbol r$。在每一小段上，元功为
\[
    \mathrm dW=\boldsymbol F\cdot\mathrm d\boldsymbol r .
\]
从起点 $a$ 到终点 $b$ 的总功为
\[
    \boxed{
    W=\int_a^b\boldsymbol F\cdot\mathrm d\boldsymbol r
    } .
\]
在直角坐标系中可写为
\[
    W=\int_a^b
    \left(
    F_x\,\mathrm dx+F_y\,\mathrm dy+F_z\,\mathrm dz
    \right).
\]
计算变力做功时，应沿实际路径积分，不能随意取始末力的平均值乘总位移。

\subsection{合力做功}

若物体受到多个力
\[
    \boldsymbol F_1,\boldsymbol F_2,\cdots,\boldsymbol F_n,
\]
则合力为
\[
    \boldsymbol F_{\mathrm{net}}=\sum_{i=1}^{n}\boldsymbol F_i .
\]
合力做功为
\[
    W_{\mathrm{net}}
    =
    \int \boldsymbol F_{\mathrm{net}}\cdot\mathrm d\boldsymbol r
    =
    \sum_{i=1}^{n}
    \int \boldsymbol F_i\cdot\mathrm d\boldsymbol r
    =
    \sum_{i=1}^{n}W_i .
\]
所以，合力做功等于各分力做功的代数和。

\subsection{功率}

平均功率定义为
\[
    \bar P=\frac{\Delta W}{\Delta t}.
\]
瞬时功率为
\[
    P=\frac{\mathrm dW}{\mathrm dt}.
\]
由 $\mathrm dW=\boldsymbol F\cdot\mathrm d\boldsymbol r$ 得
\[
    \boxed{
    P=\boldsymbol F\cdot\boldsymbol v
    } .
\]
若力与速度夹角为 $\theta$，则
\[
    P=Fv\cos\theta .
\]
功率单位为瓦特：
\[
    1\,\mathrm W=1\,\mathrm{J/s}.
\]

\subsection{功的相对性}

功与参考系有关。因为位移 $\mathrm d\boldsymbol r$ 和速度 $\boldsymbol v$ 都依赖参考系，所以同一个力在不同惯性系中做功和功率一般不同。与此同时，动能也依赖参考系；只要在同一参考系内计算，功与动能之间仍满足动能定理。

\subsection{例题：变力沿直线做功}

质点沿 $x$ 轴从 $x=0$ 运动到 $x=a$。作用在质点上的力为
\[
    F_x=kx,\qquad F_y=F_z=0,
\]
其中 $k$ 为常量。求该力做的功。

\textbf{解析：} 因为质点沿 $x$ 轴运动，
\[
    \mathrm d\boldsymbol r=\mathrm dx\,\boldsymbol e_x.
\]
元功为
\[
    \mathrm dW=\boldsymbol F\cdot\mathrm d\boldsymbol r=kx\,\mathrm dx.
\]
故
\[
    W=\int_0^a kx\,\mathrm dx=\frac{1}{2}ka^2.
\]
因此
\[
    \boxed{W=\frac{1}{2}ka^2}.
\]
本题体现了变力做功必须积分，不能简单套用恒力做功公式。

\subsection{常用公式}

\[
    W=\boldsymbol F\cdot\boldsymbol s=Fs\cos\theta
\]
\[
    \boxed{
    W=\int_a^b\boldsymbol F\cdot\mathrm d\boldsymbol r
    =
    \int_a^b
    \left(
    F_x\,\mathrm dx+F_y\,\mathrm dy+F_z\,\mathrm dz
    \right)
    }
\]
\[
    \boxed{
    P=\frac{\mathrm dW}{\mathrm dt}=\boldsymbol F\cdot\boldsymbol v
    }
\]

\newpage

\section{动能定理}

动能定理将合外力做功与质点动能变化联系起来。它是牛顿第二定律的能量形式，在处理已知位移、求速率的问题时尤其方便。

\subsection{质点动能}

质点动能定义为
\[
    \boxed{
    E_k=\frac{1}{2}mv^2
    } .
\]
动能是标量，只与速率大小有关，与速度方向无关。由于速度依赖参考系，动能也依赖参考系。

\subsection{动能定理的推导}

对质量恒定的质点 $m$，在惯性参考系中，在合外力 $\boldsymbol F_{\mathrm{net}}$ 作用下运动。由牛顿第二定律
\[
    \boldsymbol F_{\mathrm{net}}=m\boldsymbol a
\]
可得合外力元功
\[
    \mathrm dW_{\mathrm{net}}
    =
    \boldsymbol F_{\mathrm{net}}\cdot\mathrm d\boldsymbol r
    =
    m\boldsymbol a\cdot\mathrm d\boldsymbol r .
\]
又因为
\[
    \boldsymbol a=\frac{\mathrm d\boldsymbol v}{\mathrm dt},
    \qquad
    \mathrm d\boldsymbol r=\boldsymbol v\,\mathrm dt,
\]
所以
\[
    \mathrm dW_{\mathrm{net}}
    =
    m\frac{\mathrm d\boldsymbol v}{\mathrm dt}\cdot\boldsymbol v\,\mathrm dt
    =
    m\boldsymbol v\cdot\mathrm d\boldsymbol v .
\]
由
\[
    \mathrm d(v^2)
    =
    \mathrm d(\boldsymbol v\cdot\boldsymbol v)
    =
    2\boldsymbol v\cdot\mathrm d\boldsymbol v
\]
可得
\[
    m\boldsymbol v\cdot\mathrm d\boldsymbol v
    =
    \mathrm d\left(\frac{1}{2}mv^2\right).
\]
因此
\[
    \mathrm dW_{\mathrm{net}}=\mathrm dE_k.
\]
从初态到末态积分，得到
\[
    \boxed{
    W_{\mathrm{net}}
    =
    \Delta E_k
    =
    \frac{1}{2}mv_2^2-\frac{1}{2}mv_1^2
    } .
\]

\subsection{使用步骤}

应用动能定理时，可按以下顺序处理：

\begin{enumerate}
    \item 选定研究对象；
    \item 明确初态和末态；
    \item 判断哪些力做功，哪些力不做功；
    \item 写出合外力做功 $W_{\mathrm{net}}$；
    \item 写出动能变化 $\Delta E_k$；
    \item 建立方程
    \[
        W_{\mathrm{net}}=\Delta E_k
    \]
    并求解未知量。
\end{enumerate}

动能定理不直接涉及时间，因此当题目只关心经过某段位移后的速度大小时，通常比逐步求加速度和时间更简洁。

\subsection{与动量定理的区别}

动量定理为
\[
    \boldsymbol I=\Delta\boldsymbol p,
\qquad
    \boldsymbol I=\int\boldsymbol F\,\mathrm dt,
\]
它体现力对时间的累积。动能定理为
\[
    W_{\mathrm{net}}=\Delta E_k,
\qquad
    W=\int\boldsymbol F\cdot\mathrm d\boldsymbol r,
\]
它体现力对空间位移的累积。

因此：

\begin{enumerate}
    \item 冲量是矢量，功是标量；
    \item 动量变化由冲量决定，动能变化由合外力做功决定；
    \item 动量定理保留速度方向信息，动能定理主要处理速率大小变化。
\end{enumerate}

\subsection{例题：粗糙水平面上的制动距离}

质量为 $m$ 的物块以初速度 $v_0$ 在粗糙水平面上滑行，动摩擦因数为 $\mu$。求物块停止前滑行距离 $s$。

\textbf{解析：} 物块受到重力、支持力和滑动摩擦力。重力和支持力都垂直于位移，不做功；滑动摩擦力方向与位移方向相反，大小为
\[
    f_k=\mu mg.
\]
摩擦力做功为
\[
    W_f=-f_ks=-\mu mgs.
\]
物块末速度为零，所以
\[
    \Delta E_k=0-\frac{1}{2}mv_0^2.
\]
由动能定理
\[
    -\mu mgs=-\frac{1}{2}mv_0^2.
\]
解得
\[
    \boxed{
    s=\frac{v_0^2}{2\mu g}
    } .
\]

\subsection{例题：斜向拉力做功}

质量为 $m$ 的物块在水平面上运动，受到大小为 $F$、与水平面夹角为 $\theta$ 的拉力作用，水平位移为 $s$。若忽略摩擦，求拉力做功与物块动能变化。

\textbf{解析：} 位移方向水平，拉力沿位移方向的分量为
\[
    F_x=F\cos\theta.
\]
拉力做功为
\[
    W=Fs\cos\theta.
\]
重力和支持力都垂直于水平位移，不做功。由动能定理
\[
    \Delta E_k=W=Fs\cos\theta.
\]
因此
\[
    \boxed{
    \Delta E_k=Fs\cos\theta
    } .
\]
这说明并非整个拉力都改变动能，只有沿位移方向的分量做功。

\subsection{常用公式}

\[
    \boxed{
    E_k=\frac{1}{2}mv^2
    }
\]
\[
    \boxed{
    W_{\mathrm{net}}=\Delta E_k
    =
    \frac{1}{2}mv_2^2-\frac{1}{2}mv_1^2
    }
\]
\[
    \boldsymbol I=\int\boldsymbol F\,\mathrm dt,\qquad
    W=\int\boldsymbol F\cdot\mathrm d\boldsymbol r.
\]

